%قبل از شروع مسئله ترجیحاً چیزی ننویسید. اما اگر خواستید بنویسید، در قالب لیستی (فقرات و شمارش) بنویسید.
%برای اضافه کردن نمره، بعدی از شروع‌مسئله، [x] را اضافه کنید و به جای x  نمره مورد نظر را قرار دهید. برای زیرمسئله‌ها هم به همین صورت است.
%\شروع‌مسئله[10]
\شروع‌مسئله
% صورت اصلی مسئله را در این‌جا بنویسید.

 فرض کنید $n$ عدد به صورت لیست پیوندی به شما داده شده است. با کمک گرفتن از مرتب‌سازی ادغامی و پیچیدگی زمانی
  $O(n \log{n})$، اعداد را در شرایط زیر مرتب کنید.

%اگر مسئله چند قسمت داشت به این صورت عمل کنید:
\شروع{زیرمسائل}
\زیرمسئله
استفاده از حافظه اضافی از مرتبه
 $O(\log{n})$
\زیرمسئله
استفاده از حافظه اضافی از مرتبه
$O(1)$
\footnote{
دقت کنید که استفاده از رابطه بازگشتی که تا عمق
 $h$ 
 پایین می‌رود، به اندازه
  $O(h)$
   حافظه مصرف می‌کند.
   همچنین این سوال به صورت اضافی مطرح شده است و در کوئیز نمی‌آید. اما حل آن توصیه می‌شود.
}
%\زیرمسئله % مسئلۀ ۱ت
\پایان{زیرمسائل}

\solution{
	
	با استفاده از سه تابع زیر می‌توان لیست پیوندی را مرتب کرد :
	\lr{
		\begin{algorithm}[H]
			\begin{algorithmic}
				\State $ head \leftarrow fisrt \; element \; of \;  linked \; list$
				\State $ *head \leftarrow pointer \; to\; head $
				\Procedure{$merge\_sort(*head)$}{}
				\If{$head = null \;||\; head\rightarrow next = null$}
				\State $return$
				\EndIf
				\State $Split(head, \&first\_half, \&second\_half)$
				\State $merge\_sort(\&first\_half)$
				\State $merge\_sort(\&second\_half)$
				\State $*head = merge(\&first\_half, \&second\_half)$
				\EndProcedure
				\Procedure{$split(*head, *first\_half, *second\_half)$}{}
				\State $slow = head$
				\State $fast = head\rightarrow next$
				\While {$fast \neq null$}
				\State $fast = fast\rightarrow next$
				\If{$fast \neq null$}
				\State $slow = slow\rightarrow next$
				\State $fast = fast\rightarrow next$
				\EndIf
				\EndWhile 
				\State $*first\_half = head$
				\State $*second\_half = slow\rightarrow next$
				\EndProcedure
				\Procedure{$merge(*first\_half, *second\_half)$}{}
				\If{$first\_half = null$}
				\State $return \; second\_half$
				\EndIf
				\If{$second\_half = null$}
				\State $return \; first\_half$
				\EndIf
				\If{$first\_half \leq second\_half$}
				\State $result = first\_half$
				\State $result \rightarrow next = merge(first\_half, second\_half \rightarrow next)$
				\Else 
				\State $result = second\_half$
				\State $result \rightarrow next = merge(first\_half, second\_half\rightarrow next,)$
				\EndIf
				\State $return \; result$
				\EndProcedure
			\end{algorithmic}
		\end{algorithm}
	}

قسمت ۲:

\hyperlink{https://www.chiark.greenend.org.uk/~sgtatham/algorithms/listsort.html}{لینک حل}
\begin{latin}
	\lstinputlisting[language=C++]{src/2-b.cpp}
\end{latin}


}



%%%%%%%%%%%%%%%%%%%%%%%%%%%%%%%%%%%%%%%%%%%%%%%%%%%%%
% برخی دستورات مفید برای تمرین‌ها. %
%%%%%%%%%%%%%%%%%%%%%%%%%%%%%%%%%%%%%%%%%%%%%%%%%%%%%

% فرمول‌های نمونه:
%$f(x) = \Theta(x)$
%$T(x, y) \leq \log(x) + 2^y + \binom{2n}{n}$
% {\tt code\_function}


% لیست موارد بدون شماره‌گذاری:
%\شروع{فقرات}
%    \فقره مورد اول 
%        \شروع{فقرات}
%            \فقره مورد اول داخلی 
%                \شروع{فقرات}
%                    \فقره اولین مورد خیلی داخلی 
%                    \فقره دومین مورد خیلی داخلی 
%                \پایان{فقرات}
%            \فقره مورد دوم داخلی 
%        \پایان{فقرات}
%    \فقره مورد دوم 
%\پایان{فقرات}

% لیست موارد شماره‌گذاری شده:
%\شروع{شمارش}
%    \فقره یک مورد 
%    \فقره یک مورد دیگر 
%    \فقره باز هم یک مورد دیگر
%\پایان{شمارش}

% برای نوشتن یک فرمول چند خطی که یک شماره بخوره. برای حالت‌های دیگر به سایت زیر مراجعه کنید:
% https://www.overleaf.com/learn/latex/Aligning_equations_with_amsmath
%\begin{align} 
%    \begin{split}
%        (x+y)^3 &= (x+y)^2(x+y) \\
%                &= (x^2+2xy+y^2)(x+y) \\
%                &= (x^3+2x^2y+xy^2) + (x^2y+2xy^2+y^3) \\
%                &= x^3+3x^2y+3xy^2+y^3
%    \end{split}                                 
%\end{align}

% نوشتن ماتریس و آرایه دوبعدی
%\begin{align}
%    A = 
%    \begin{bmatrix}
%        A_{11} & A_{21} \\
%        A_{21} & A_{22}
%    \end{bmatrix}
%\end{align}

% اضافه کردن تصویر با شرح. برای اضافه کردن تصویر بدون شرح، فقط دستور درج‌تصویر را استفاده کنید.
%\شروع{شکل}[h!]
%\تنظیم‌ازوسط
%  \درج‌تصویر[پهنا=0.5\پهنای‌متن]{img.jpg}
%\شرح{شرح تصویر}
%\برچسب{تصویر:۱}
%\پایان{شکل}
%\پایان{زیرمسائل}

%کد زیر را ببینید.
%%\lstinputlisting{source_filename.py}
%\begin{latin}
%\begin{lstlisting}
%# Import the modules
%import sys
%import random
%
%ans = True
%
%while ans:
%question = raw_input("Ask the magic 8 ball a question: (press enter to quit) ")
%
%answers = random.randint(1,8)
%
%if question == "":
%sys.exit()
%
%elif answers == 1:
%print "It is certain"
%
%elif answers == 2:
%print "Outlook good"
%
%elif answers == 3:
%print "You may rely on it"
%
%elif answers == 4:
%print "Ask again later"
%
%elif answers == 5:
%print "Concentrate and ask again"
%
%elif answers == 6:
%print "Reply hazy, try again"
%
%elif answers == 7:
%print "My reply is no"
%
%elif answers == 8:
%print "My sources say no"
%\end{lstlisting}
%\end{latin}
